%==============================================================
% ΚΕΦΑΛΑΙΟ 1 — ΕΙΣΑΓΩΓΗ
%==============================================================
\chapter{Εισαγωγή}
\label{ch:intro}

%--------------------------------------------------------------
\section{Πλαίσιο και Κίνητρο}
\label{sec:motivation}
%--------------------------------------------------------------

Η αξιολόγηση του πιστωτικού κινδύνου αποτελεί μία από τις
θεμελιώδεις λειτουργίες κάθε χρηματοπιστωτικού ιδρύματος. Από αυτή
την αξιολόγηση εξαρτάται η απόφαση χορήγησης δανείου, ο
καθορισμός του επιτοκίου και τελικά η κερδοφορία και η σταθερότητα
του ίδιου του ιδρύματος \citep{basel2006}. Η ιστορική πορεία των
μεθόδων που χρησιμοποιούνται για τον σκοπό αυτό αντικατοπτρίζει την
εξέλιξη της ίδιας της επιστήμης δεδομένων: από την υποκειμενική
κρίση εμπείρων αναλυτών, στα στατιστικά scorecards και τη Logistic
Regression και πλέον στα σύγχρονα μοντέλα μηχανικής μάθησης.

Τα μοντέλα μηχανικής μάθησης παρέχουν συστηματικά υψηλότερη
προβλεπτική ακρίβεια έναντι των παραδοσιακών εργαλείων
\citep{hastie2009}, γεγονός που εξηγεί τη σταδιακή υιοθέτησή τους
από τον τραπεζικό κλάδο. Η υιοθέτηση αυτή όμως δεν είναι χωρίς
κόστος. Όσο πιο πολύπλοκο γίνεται ένα μοντέλο, τόσο πιο δύσκολη
καθίσταται η κατανόηση των αποφάσεών του και η εξήγησή τους τόσο
στους δανειολήπτες όσο και στις εποπτικές αρχές. Η ένταση μεταξύ
απόδοσης και διαφάνειας τοποθετεί το πεδίο της επικύρωσης μοντέλων
στο κέντρο της σύγχρονης συζήτησης για την αξιόπιστη χρήση της
τεχνητής νοημοσύνης στις χρηματοπιστωτικές υπηρεσίες.

Παράλληλα, το κανονιστικό πλαίσιο γίνεται διαρκώς πιο απαιτητικό. Οι
κατευθύνσεις της Επιτροπής Βασιλείας (Basel~II/III), το SR Letter
11-7 της Federal Reserve και οι EBA Guidelines επιβάλλουν αυστηρές
απαιτήσεις τεκμηρίωσης, ερμηνείας και παρακολούθησης κάθε εσωτερικού
μοντέλου που χρησιμοποιείται για τον υπολογισμό κεφαλαιακών
απαιτήσεων \citep{basel2017, sr117, eba2023}. Η συμμόρφωση με αυτές
τις απαιτήσεις δεν είναι προαιρετική και δεν εξαντλείται στην
αξιολόγηση της απόδοσης — εκτείνεται σε άξονες όπως η ερμηνευσιμότητα,
η ανίχνευση μεροληψίας και η παρακολούθηση μετά το deployment.

Η συνύπαρξη των δύο αυτών δυνάμεων ---της τεχνολογικής εξέλιξης που
ωθεί προς πιο σύνθετα μοντέλα και του κανονιστικού πλαισίου που
απαιτεί διαρκώς υψηλότερη διαφάνεια--- δημιουργεί την ανάγκη για
συστηματικά πλαίσια επικύρωσης. Η παρούσα εργασία τοποθετείται
ακριβώς σε αυτό το σημείο τομής.

%--------------------------------------------------------------
\section{Στόχοι της Εργασίας}
\label{sec:objectives}
%--------------------------------------------------------------

Η εργασία θέτει τέσσερις συγκεκριμένους στόχους.

\paragraph{Πρώτος στόχος} είναι η συστηματική σύγκριση τριών
αντιπροσωπευτικών μοντέλων διαφορετικής πολυπλοκότητας: της
Logistic Regression ως παραδοσιακού σημείου αναφοράς, του XGBoost
ως κύριου μοντέλου ανάλυσης και ενός Feed-Forward Neural Network ως
benchmark αυξημένης πολυπλοκότητας. Η σύγκριση εστιάζει όχι μόνο
στην προβλεπτική απόδοση αλλά και στις απαιτήσεις επικύρωσης που
συνοδεύουν κάθε μοντέλο.

\paragraph{Δεύτερος στόχος} είναι η αξιολόγηση των μοντέλων σε
τέσσερις άξονες παράλληλα: απόδοση, ερμηνευσιμότητα, μεροληψία και
παρακολούθηση. Η ταυτόχρονη ανάλυση όλων των αξόνων επιτρέπει
συμπεράσματα που δεν θα προέκυπταν από την εξέταση ενός μόνο εξ
αυτών.

\paragraph{Τρίτος στόχος} είναι η ποσοτική αποτύπωση του κόστους
που συνεπάγεται η αυξημένη πολυπλοκότητα ενός μοντέλου τόσο σε
υπολογιστικούς πόρους όσο και σε ευχέρεια ερμηνείας και διαχείρισης
μεροληψίας. Στόχος είναι να καταστεί σαφές πού η πολυπλοκότητα
δικαιολογείται και πού όχι.

\paragraph{Τέταρτος στόχος} είναι η σύνθεση των ευρημάτων σε ένα
συγκεκριμένο πλαίσιο επικύρωσης που να μπορεί να αποτελέσει
πρακτικό εργαλείο. Το πλαίσιο αυτό περιλαμβάνει αρχές σχεδιασμού,
πρότυπη ροή εργασιών και διαφοροποιημένες απαιτήσεις ανά επίπεδο
πολυπλοκότητας.

%--------------------------------------------------------------
\section{Ερευνητικό Ερώτημα}
\label{sec:research_questions}
%--------------------------------------------------------------

Το κεντρικό ερευνητικό ερώτημα της εργασίας διατυπώνεται ως εξής:

\begin{quote}
\textit{Πώς επικυρώνουμε, ερμηνεύουμε και παρακολουθούμε μοντέλα
μηχανικής μάθησης στον πιστωτικό κίνδυνο, και πώς μεταβάλλονται οι
απαιτήσεις αυτές καθώς αυξάνεται η πολυπλοκότητα του μοντέλου;}
\end{quote}

Η διατύπωση αυτή είναι σκόπιμα διπλή. Το πρώτο σκέλος αφορά τη
\emph{μεθοδολογία} επικύρωσης, ερμηνείας και παρακολούθησης. Το
δεύτερο σκέλος αφορά τη \emph{σύγκριση} των απαιτήσεων αυτών μεταξύ
μοντέλων διαφορετικής πολυπλοκότητας. Η απάντηση στο ερώτημα αυτό
δεν είναι απλώς μια καταγραφή τεχνικών, αλλά μια ποιοτική κρίση για
το ποιες από τις απαιτήσεις κλιμακώνονται γραμμικά με την
πολυπλοκότητα και ποιες όχι. Στο συμπέρασμα αυτό βασίζεται και η
πρόταση πλαισίου επικύρωσης που παρουσιάζεται στο
Κεφάλαιο~\ref{ch:framework}.

Για τη διερεύνηση του ερωτήματος αυτού χρησιμοποιείται η βάση
δεδομένων FICO HELOC, η οποία αποτελεί ευρέως εδραιωμένη επιλογή
στη βιβλιογραφία της ερμηνεύσιμης μηχανικής μάθησης στον πιστωτικό
κίνδυνο \citep{fico2018}.

%--------------------------------------------------------------
\section{Δομή της Εργασίας}
\label{sec:structure}
%--------------------------------------------------------------

Η εργασία αναπτύσσεται σε έξι κεφάλαια. Στο
Κεφάλαιο~\ref{ch:theory} παρουσιάζεται το θεωρητικό και κανονιστικό
πλαίσιο, με ανάλυση της έννοιας του πιστωτικού κινδύνου, της
ιστορικής εξέλιξης των μεθόδων credit scoring, των μοντέλων
μηχανικής μάθησης που χρησιμοποιούνται και των τεχνικών επικύρωσης,
ερμηνείας, ανίχνευσης μεροληψίας και παρακολούθησης.

Στο Κεφάλαιο~\ref{ch:methodology} περιγράφονται αναλυτικά η βάση
δεδομένων, η προεπεξεργασία της, οι παράμετροι των τριών μοντέλων
και η μεθοδολογία επικύρωσης που εφαρμόζεται.

Στο Κεφάλαιο~\ref{ch:results} παρουσιάζονται τα πειραματικά
αποτελέσματα οργανωμένα στους τέσσερις άξονες αξιολόγησης
(απόδοση, ερμηνευσιμότητα, μεροληψία, παρακολούθηση), με ιδιαίτερη
έμφαση στην εξειδικευμένη ανάλυση του XGBoost ως κύριου μοντέλου
της εργασίας.

Στο Κεφάλαιο~\ref{ch:framework} συντίθενται τα ευρήματα σε μια
συγκεκριμένη πρόταση πλαισίου επικύρωσης, η οποία περιλαμβάνει
αρχές σχεδιασμού, πρότυπη ροή εργασιών και διαφοροποιημένες
απαιτήσεις ανά επίπεδο πολυπλοκότητας.

Τέλος, στο Κεφάλαιο~\ref{ch:conclusions} συνοψίζονται τα κύρια
συμπεράσματα, διατυπώνεται η ευθεία απάντηση στο ερευνητικό ερώτημα
και προτείνονται κατευθύνσεις για περαιτέρω έρευνα.